\documentclass[11pt,a4paper]{article}
\usepackage[utf8]{inputenc}
\usepackage[spanish]{babel}
\usepackage{amsmath, amssymb}
\usepackage{geometry}
\geometry{top=2.5cm, bottom=2.5cm, left=2.5cm, right=2.5cm}

\begin{document}

\section*{Piense, dialogue y comparta}

\begin{itemize}
    \item[\textbf{1.}] Usted observa a su pequeño sobrino rodar canicas desde la parte superior de una escalera. Esta tiene 12 escalones, cada uno de $30.0 \text{ cm}$ de profundidad horizontal, separados por $20.0 \text{ cm}$ verticalmente. Las canicas salen horizontalmente del escalón superior y se proyectan en el aire, rebotando en los escalones hasta que llegan a la planta baja. Esto le hace preguntarse lo siguiente: (a) ¿Qué tan rápido debe rodar la canica para que no rebote en el \textit{primer} escalón debajo del piso superior? (b) ¿Qué tan rápido debe rodar la canica para que no rebote en el \textit{segundo} escalón debajo del piso superior? (c) ¿Es posible que su pequeño sobrino ruede la canica lo suficientemente rápido como para no pegar en \textit{todos} los escalones? (d) Suponga que la canica se proyecta con una rapidez tal que llega al sexto escalón y rebota hacia arriba con el mismo ángulo con el que pegó en el escalón, con la misma rapidez. Argumente por qué la canica no pegará en otro escalón antes de pegar en el piso de la planta baja.
    \vspace*{9cm}

    \item[\textbf{2.}] \textbf{ACTIVIDAD} Coloque una moneda en la esquina de una mesa como se muestra en la vista superior de la figura TP4.2. Coloque una regla junto a la moneda y otra moneda en la parte superior de la regla que cuelga del borde de la mesa. Sostenga el extremo de la regla sobre la mesa con una mano y utilice la otra mano para mover el extremo de la regla con la moneda paralela a la superficie de la mesa. Esto proyectará la moneda que está en la esquina de la mesa en una dirección horizontal. Al mismo tiempo, la regla se deslizará hacia afuera por debajo de la segunda moneda, que caerá hacia abajo. Use el grabador de audio de su teléfono inteligente, haga una grabación de audio de las dos monedas cayendo. A partir de la grabación, determine el intervalo de tiempo entre la llegada de las dos monedas al suelo. ¿Cuál debe ser teóricamente el intervalo de tiempo?
    \vspace*{8cm}
\end{itemize}

\newpage

\section*{Problemas}

\subsection*{SECCIÓN 4.1 Vectores de posición, velocidad y aceleración}

\begin{itemize}
    \item[\textbf{1.}] Supongamos que el vector de posición de una partícula está dado como una función del tiempo por $\vec{r}(t) = x(t)\hat{i} + y(t)\hat{j}$, con $x(t) = at + b$ y $y(t) = ct^2 + d$, donde $a = 1.00 \text{ m/s}$, $b = 1.00 \text{ m}$, $c = 0.125 \text{ m/s}^2$ y $d = 1.00 \text{ m}$. (a) Calcule la velocidad promedio durante el intervalo de tiempo de $t = 2.00 \text{ s}$ a $t = 4.00 \text{ s}$. (b) Determine la velocidad y la rapidez en $t = 2.00 \text{ s}$.
    \vspace*{6cm}

    \item[\textbf{2.}] Las coordenadas de un objeto que se mueve en el plano $xy$ varían con el tiempo de acuerdo con $x = -(5.00) \operatorname{sen}(\omega t)$ y $y = (4.00) - (5.00) \cos(\omega t)$, donde $\omega$ es una constante, $x$ y $y$ están en metros y $t$ está en segundos. (a) Determine las componentes de velocidad del objeto en $t = 0$. (b) Determine las componentes de la aceleración del objeto en $t = 0$. (c) Escriba expresiones para el vector de posición, el vector velocidad y el vector aceleración del objeto en cualquier tiempo $t > 0$. (d) Describa la trayectoria del objeto en una gráfica $xy$.
    \vspace*{8cm}
\end{itemize}

\newpage

\subsection*{SECCIÓN 4.2 Movimiento en dos dimensiones con aceleración constante}

\begin{itemize}
    \item[\textbf{3.}] El vector de posición de una partícula varía en el tiempo de acuerdo con la expresión $\vec{r} = 3.00\hat{i} - 6.00t^2\hat{j}$, donde $\vec{r}$ está en metros y $t$ en segundos. (a) Encuentre una expresión para la velocidad de la partícula como función del tiempo. (b) Determine la aceleración de la partícula como función del tiempo. (c) Calcule la posición y velocidad de la partícula en $t = 1.00 \text{ s}$.
    \vspace*{6cm}

    \item[\textbf{4.}] No es posible ver objetos muy pequeños, como virus, utilizando un microscopio de luz ordinario. Sin embargo, en un microscopio electrónico sí es posible verlos ya que este utiliza un haz de electrones en lugar de un haz de luz. La microscopía electrónica ha demostrado ser invaluable para la investigación de virus, membranas celulares y estructuras subcelulares, superficies bacterianas, receptores visuales, cloroplastos y propiedades contráctiles de los músculos. Las ``lentes'' de un microscopio electrónico consisten en campos eléctricos y magnéticos que controlan el haz de electrones. Como ejemplo de la manipulación de un haz de electrones, considere un electrón que viaja lejos del origen a lo largo del eje $x$ en el plano $xy$ con velocidad inicial $\vec{v}_i = v_i\hat{i}$. Al pasar por la región de $x = 0$ a $x = d$, el electrón experimenta una aceleración $\vec{a} = a_x\hat{i} + a_y\hat{j}$, donde $a_x$ y $a_y$ son constantes. Para el caso $v_i = 1.80 \times 10^7 \text{ m/s}$, $a_x = 8.00 \times 10^{14} \text{ m/s}^2$ y $a_y = 1.60 \times 10^{15} \text{ m/s}^2$, determine en $x = d = 0.010\ 0 \text{ m}$ (a) la posición del electrón, (b) la velocidad del electrón, (c) la rapidez del electrón y (d) la dirección de desplazamiento del electrón (es decir, el ángulo entre su velocidad y el eje de $x$).
    \vspace*{8cm}

    \item[\textbf{5.}] \textbf{Problema de repaso.} Una moto de nieve está originalmente en el punto con el vector de posición $29.0 \text{ m}$ a $95.0^\circ$ en sentido contrario a las manecillas del reloj desde el eje $x$, se mueve con una velocidad de $4.50 \text{ m/s}$ a $40.0^\circ$ y con aceleración constante de $1.90 \text{ m/s}^2$ a $200^\circ$. Después de que han transcurrido $5.00 \text{ s}$, encuentre (a) su velocidad y (b) su vector de posición.
    \vspace*{6cm}
\end{itemize}

\newpage

\subsection*{SECCIÓN 4.3 Movimiento de un proyectil}

\begin{itemize}
    \item[\textbf{6.}] En un bar local, un cliente desliza sobre la barra un tarro de cerveza vacío para que lo vuelvan a llenar. La altura de la barra es $h$. El cantinero está momentáneamente distraído y no ve el tarro, que se desliza de la barra y golpea el suelo a $d$ de la base de la barra. (a) ¿Con qué velocidad el tarro dejó la barra? (b) ¿Cuál fue la dirección de la velocidad del tarro justo antes de golpear el suelo?
    \vspace*{6cm}

    \item[\textbf{7.}] Los reyes mayas y muchos equipos de deportes de la escuela se llaman puma, pantera o león de montaña —el \textit{Felis concolor}—, que es el mejor saltador entre los animales. Puede saltar a una altura de $12.0 \text{ pies}$ desde el suelo con un ángulo de $45.0^\circ$. ¿Con qué rapidez, en unidades SI, se despega del suelo para dar este salto?
    \vspace*{5cm}

    \item[\textbf{8.}] Un proyectil se dispara en tal forma que su alcance horizontal es igual a tres veces su altura máxima. ¿Cuál es el ángulo de proyección?
    \vspace*{5cm}

    \item[\textbf{9.}] La rapidez de un proyectil cuando alcanza su altura máxima es la mitad de su velocidad cuando está en la mitad de su altura máxima. ¿Cuál es el ángulo de proyección inicial del proyectil?
    \vspace*{6cm}
\end{itemize}

\newpage

\begin{itemize}
    \item[\textbf{10.}] Una roca se lanza hacia arriba desde el suelo de forma que la altura máxima de su vuelo es igual a su alcance horizontal $R$. (a) ¿A qué ángulo $\theta$ se lanza la roca? (b) En términos del alcance original $R$, ¿cuál es el alcance $R_{\text{máx}}$ que puede lograr la roca si se lanza a la misma rapidez, pero en ángulo óptimo para el alcance máximo? (c) \textbf{¿Qué pasaría si?} ¿Su respuesta al inciso (a) cambiaría si la roca se lanza con la misma rapidez en un planeta diferente? Explique.
    \vspace*{6cm}

    \item[\textbf{11.}] Un bombero, a una distancia $d$ de un edificio en llamas, dirige un chorro de agua desde una manguera en un ángulo $\theta_i$ sobre la horizontal, como se muestra en la figura P4.11. Si la rapidez inicial del chorro es $v_i$, ¿en qué altura $h$ el agua golpea al edificio?
    \vspace*{6cm}

    \item[\textbf{12.}] Una estrella de basquetbol cubre horizontalmente $2.80 \text{ m}$ en un salto para encestar la pelota (figura P4.12a). Su movimiento a través del espacio se representa igual que el de una partícula en su \textit{centro de masa}, que se definirá en el capítulo 9. Su centro de masa está a una altura de $1.02 \text{ m}$ cuando deja el suelo. Llega altura máxima de $1.85 \text{ m}$ sobre el suelo y está a una elevación de $0.900 \text{ m}$ cuando toca el suelo de nuevo. Determine: (a) su tiempo de vuelo (su ``tiempo colgado''), (b) sus componentes de velocidad horizontal y (c) vertical en el instante de despegar, y (d) su ángulo de despegue. (e) En comparación, determine el tiempo colgado de un ciervo cola blanca que da un salto (figura P4.12b) con elevaciones de centro de masa $y_i = 1.20 \text{ m}$, $y_{\text{máx}} = 2.50 \text{ m}$ y $y_f = 0.700 \text{ m}$.
    \vspace*{8cm}
\end{itemize}

\newpage

\begin{itemize}
    \item[\textbf{13.}] Un estudiante se encuentra en el borde de un acantilado y lanza una piedra horizontalmente sobre el borde con una velocidad de $v_i = 18.0 \text{ m/s}$. El acantilado tiene $h = 50.0 \text{ m}$ arriba del agua como se muestra en la figura P4.13. (a) ¿Cuáles son las coordenadas de la posición inicial de la piedra? (b) ¿Cuáles son los componentes de la velocidad inicial de la piedra? (c) ¿Cuál es el modelo de análisis adecuado para el movimiento vertical de la piedra? (d) ¿Cuál es el modelo de análisis adecuado para el movimiento horizontal de la piedra? (e) Escriba ecuaciones simbólicas para las componentes $x$ y $y$ de la velocidad de la piedra como una función del tiempo. (f) Escriba ecuaciones simbólicas para la posición de la piedra como una función del tiempo. (g) ¿Cuánto tiempo después de ser lanzada la piedra golpeó el agua debajo del acantilado? (h) ¿Con qué rapidez y ángulo de impacto pegó la piedra?
    \vspace*{9cm}

    \item[\textbf{14.}] La distancia récord en el deporte de lanzar discos es de $81.1 \text{ m}$. Este lanzamiento de discos fue establecido por Steve Urner de Estados Unidos en 1981. Suponiendo que el ángulo de lanzamiento inicial fue de $45^\circ$ y despreciando la resistencia del aire, determine (a) la rapidez inicial del proyectil y (b) el intervalo de tiempo total que voló el proyectil. (c) ¿Cómo cambiaría las respuestas si el alcance fuera el mismo pero el ángulo de lanzamiento fuera mayor que $45^\circ$? Explique.
    \vspace*{6cm}
\end{itemize}

\newpage

\begin{itemize}
    \item[\textbf{15.}] Se hace un jonrón al pegarle a la pelota de béisbol de manera que libre una pared de $21.0 \text{ m}$ de alto, situada a $130 \text{ m}$ de la base de \textit{home}. La pelota es golpeada con un ángulo de $35.0^\circ$ con la horizontal, y se desprecia la resistencia del aire. Encuentre (a) la rapidez inicial de la pelota, (b) el tiempo que tarda esta en llegar a la pared y (c) las componentes de su velocidad y rapidez, cuando llega a la pared. (Suponga que la pelota es golpeada a una altura de $1.00 \text{ m}$ del suelo.)
    \vspace*{6cm}

    \item[\textbf{16.}] Un proyectil es disparado desde lo alto de un acantilado de altura $h$ sobre el océano. El proyectil se disparó a un ángulo $\theta$ por encima de la horizontal y con una velocidad inicial de $v_i$. (a) Encuentre una expresión simbólica en términos de las variables $v_i$, $g$ y $\theta$ para el tiempo en que el proyectil alcanza la altura máxima. (b) Utilizando el resultado del inciso (a), encuentre una expresión para la altura máxima $h_{\text{máx}}$ sobre el océano alcanzada por el proyectil en términos de $h$, $v_i$, $g$ y $\theta$.
    \vspace*{6cm}

    \item[\textbf{17.}] Un niño está parado en un trampolín y lanza una piedra a una piscina. La piedra es lanzada desde una altura de $2.50 \text{ m}$ sobre la superficie del agua con una velocidad de $4.00 \text{ m/s}$ en un ángulo de $60.0^\circ$ sobre la horizontal. Cuando la piedra pega en la superficie del agua, inmediatamente disminuye su rapidez a la mitad de la que tenía antes de golpear el agua y la mantiene en ella. Después de que la piedra entra en el agua, se mueve en línea recta en la dirección de la velocidad que tenía cuando la golpeó. Si la piscina tiene $3.00 \text{ m}$ de profundidad, ¿cuánto tiempo transcurre entre el momento en que se lanza la piedra y cuando esta pega en el fondo de la piscina?
    \vspace*{7cm}
\end{itemize}

\newpage

\subsection*{SECCIÓN 4.4 Modelo de análisis: Partícula en movimiento circular uniforme}

\begin{itemize}
    \item[\textbf{18.}] En el ejemplo 4.6 encontramos la aceleración centrípeta de la Tierra cuando gira alrededor del Sol. Calcule la aceleración radial de un punto en la superficie de la Tierra en el ecuador debido a la rotación de la Tierra sobre su eje.
    \vspace*{6cm}

    \item[\textbf{19.}] El astronauta que orbita la Tierra en la figura P4.19 se está preparando para acoplar con un satélite \textit{Westar VI}. El satélite está en una órbita circular $600 \text{ kilómetros}$ sobre la superficie de la Tierra, donde la aceleración de caída libre es $8.21 \text{ m/s}^2$. Tome el radio de la Tierra como $6400 \text{ km}$. Determine la rapidez del satélite y el intervalo de tiempo necesario para completar una órbita alrededor de la Tierra, que es el periodo del satélite.
    \vspace*{6cm}

    \item[\textbf{20.}] Un atleta batea una pelota, conectada al extremo de una cadena, en un círculo horizontal. El atleta es capaz de girar la pelota con una rapidez de $8.00 \text{ rev/s}$ cuando la longitud de la cadena es $0.600 \text{ m}$. Cuando la longitud aumenta a $0.900 \text{ m}$, puede rotar la pelota solo $6.00 \text{ rev/s}$. (a) ¿Cuál es la tasa de rotación que da la mayor rapidez a la pelota? (b) ¿Cuál es la aceleración centrípeta de la pelota a $8.00 \text{ rev/s}$? (c) ¿Cuál es la aceleración centrípeta en $6.00 \text{ rev/s}$?
    \vspace*{6cm}

    \item[\textbf{21.}] El atleta que se muestra en la figura P4.21 rota un disco de $1.00 \text{ kg}$ a lo largo de una trayectoria circular de $1.06 \text{ m}$ de radio. La rapidez máxima del disco es $20.0 \text{ m/s}$. Determine la magnitud de la aceleración radial máxima del disco.
    \vspace*{6cm}

    \item[\textbf{22.}] Un neumático de $0.500 \text{ m}$ de radio gira a una rapidez constante de $200 \text{ rev/min}$. Encuentre la rapidez y la aceleración de una pequeña piedra atascada en la banda de rodadura de los neumáticos (en su borde exterior).
    \vspace*{6cm}
\end{itemize}

\newpage

\subsection*{SECCIÓN 4.5 Aceleraciones tangencial y radial}

\begin{itemize}
    \item[\textbf{23.}] (a) Una partícula, que se mueve con rapidez instantánea de $3.00 \text{ m/s}$ en una trayectoria con $2.00 \text{ m}$ de radio de curvatura, ¿podría tener una aceleración de $6.00 \text{ m/s}^2$ de magnitud? (b) ¿Podría tener una aceleración $4.00 \text{ m/s}^2$? En cada caso, si la respuesta es sí, explique cómo puede ocurrir; si la respuesta es no, explique por qué.
    \vspace*{6cm}

    \item[\textbf{24.}] Una bola se balancea en un círculo vertical en el extremo de una cuerda de $1.50 \text{ m}$ de largo. Cuando la bola está a $36.9^\circ$ después del punto más bajo en su viaje hacia arriba, su aceleración total es $(-22.5\hat{i} + 20.2\hat{j}) \text{ m/s}^2$. En ese instante, (a) bosqueje un diagrama vectorial que muestre las componentes de su aceleración, (b) determine la magnitud de su aceleración radial y (c) determine la rapidez y velocidad de la bola.
    \vspace*{8cm}
\end{itemize}

\newpage

\subsection*{SECCIÓN 4.6 Velocidad y aceleración relativas}

\begin{itemize}
    \item[\textbf{25.}] Un tornillo cae desde el techo de un vagón de ferrocarril en movimiento que acelera hacia el norte en proporción de $2.50 \text{ m/s}^2$. (a) ¿Cuál es la aceleración del tornillo en relación con el vagón de ferrocarril? (b) ¿Cuál es la aceleración del tornillo en relación con la Tierra? (c) Describa la trayectoria del tornillo como la ve un observador dentro del vagón. (d) Describa la trayectoria del tornillo como la ve un observador fijo en la Tierra.
    \vspace*{7cm}

    \item[\textbf{26.}] El piloto de un avión registra que la brújula indica un rumbo hacia el oeste. La rapidez del avión respecto al aire es de $150 \text{ km/h}$. El aire se está moviendo con un viento de $30.0 \text{ km/h}$ hacia el norte. Encuentre la velocidad del avión en relación con el suelo.
    \vspace*{6cm}

    \item[\textbf{27.}] Usted está tomando lecciones de vuelo de un piloto experimentado. Ambos están en el avión, usted en el asiento del piloto. La torre de control le indica por radio, que mientras usted ha estado en el aire ha surgido un viento de $25 \text{ mi/h}$ con dirección del viento perpendicular a la pista en la que planea aterrizar. El piloto le dice que su velocidad del aire normal cuando aterrice será de $80 \text{ mi/h}$ en relación con la Tierra. Esta rapidez es relativa al aire, en la dirección en la que apunta la nariz del avión. Le pide que determine el ángulo en el que el avión debe volar de lado como ``cangrejo'', es decir, el ángulo entre la línea central del avión y la línea central de la pista que permitirá que la velocidad del avión en relación con el suelo sea paralela a la pista.
    \vspace*{6cm}

    \item[\textbf{28.}] Un automóvil viaja hacia el este con una rapidez de $50.0 \text{ km/h}$. Gotas de lluvia caen con una rapidez constante en vertical respecto de la Tierra. Las trazas de la lluvia en las ventanas laterales del automóvil forman un ángulo de $60.0^\circ$ con la vertical. Encuentre la velocidad de la lluvia en relación con (a) el automóvil y (b) la Tierra.
    \vspace*{6cm}
\end{itemize}

\newpage

\begin{itemize}
    \item[\textbf{29.}] Un estudiante de ciencias viaja en el vagón de un tren que corre a lo largo de una pista horizontal recta con una rapidez constante de $10.0 \text{ m/s}$. El estudiante lanza una bola en el aire a lo largo de una trayectoria que él juzga con un ángulo inicial de $60.0^\circ$ sobre la horizontal y está en línea con la vía. La profesora del estudiante, que está de pie en el suelo cerca de ahí, observa que la bola se eleva verticalmente. ¿Qué tan alto ve que esta se eleva?
    \vspace*{6cm}

    \item[\textbf{30.}] Un río tiene una rapidez estable de $5.00 \text{ m/s}$. Un estudiante nada corriente arriba una distancia de $1.00 \text{ km}$ y de regreso al punto de partida. (a) Si el estudiante puede nadar a una rapidez de $1.20 \text{ m/s}$ en aguas tranquilas, (a) ¿en cuánto tiempo hace el viaje redondo? (b) ¿Cuánto tempo requiere si las aguas están tranquilas? (c) Intuitivamente, ¿por qué nadar toma más tiempo cuando hay corriente?
    \vspace*{7cm}

    \item[\textbf{31.}] Un río tiene una rapidez estable $v$. Un estudiante nada corriente arriba una distancia $d$ y de regreso al punto de partida. El estudiante puede nadar a una rapidez $c$ en aguas tranquilas. (a) En términos de $d$, $v$ y $c$, ¿en cuánto tiempo hace el viaje redondo? (b) ¿Cuánto tiempo requiere si las aguas están tranquilas? (c) ¿Qué intervalo de tiempo es más largo? Explique si este siempre es más grande.
    \vspace*{7cm}

    \item[\textbf{32.}] Usted está participando en una pasantía de verano con la Guardia costera. Se le ha asignado el deber de determinar la dirección en la que una lancha de motor de la guardia costera debe viajar para interceptar naves no identificadas. Un día, el operador de radar detecta una nave no identificada a una distancia de $20.0 \text{ km}$ de la instalación del radar en la dirección $15.0^\circ$ al noreste. La nave viaja a $26.0 \text{ km/h}$ en una dirección a $40.0^\circ$ al noreste. La guardia costera desea enviar una lancha rápida, que se mueva a $50.0 \text{ km/h}$, para viajar en línea recta desde la instalación del radar para interceptar e investigar la nave, y le piden la orientación que debe tomar la lancha rápida. Exprese la dirección como una orientación de brújula respecto al norte.
    \vspace*{7cm}

    \item[\textbf{33.}] Un camión de granja se dirige al este con una velocidad constante de $9.50 \text{ m/s}$ en un tramo horizontal ilimitado del camino. Un niño se monta en la parte trasera del camión y lanza una lata de refresco hacia arriba (figura P4.33) y atrapa el proyectil en el mismo punto, pero $16.0 \text{ m}$ más lejos en el camino. (a) En el marco de referencia el camión, ¿a qué ángulo con la vertical el niño lanza la lata? (b) ¿Cuál es la rapidez inicial de la lata en relación con el camión? (c) ¿Cuál es la forma de la trayectoria de la lata como la ve el niño? Un observador en el suelo observa al niño lanzar la lata y atraparla. En este marco de referencia, (d) describa la forma de la trayectoria de la lata y (e) determine la velocidad inicial de la lata.
    \vspace*{9cm}
\end{itemize}

\newpage

\subsection*{PROBLEMAS ADICIONALES}

\begin{itemize}
    \item[\textbf{34.}] Una bola en el extremo de una cuerda gira alrededor en un círculo horizontal de radio $0.300 \text{ m}$. El plano del círculo está a $1.20 \text{ m}$ arriba del suelo. La cuerda se rompe y la bola aterriza a $2.00 \text{ m}$ (horizontalmente) alejándose del punto sobre el suelo directamente debajo de la ubicación de la bola cuando se rompe la cuerda. Encuentre la aceleración radial de la pelota durante su movimiento circular.
    \vspace*{6cm}

    \item[\textbf{35.}] ¿Por qué no es posible la siguiente situación? Un adulto de proporciones normales camina enérgicamente a lo largo de una línea recta en la dirección $+x$, parado y sosteniendo su brazo derecho vertical junto a su cuerpo para que el brazo no se balancee. Su mano derecha sostiene una pelota a su lado a una distancia $h$ arriba del piso. Cuando la pelota pasa por encima de un punto marcado como $x = 0$ en el piso horizontal, abre los dedos para soltar la pelota a partir del reposo en relación con su mano. La pelota pega en el suelo por primera vez en la posición $x = 7.00h$.
    \vspace*{6cm}

    \item[\textbf{36.}] Una partícula parte del origen con velocidad $5\hat{i} \text{ m/s}$ en $t = 0$ y se mueve en el plano $xy$ con una aceleración variable dada por $\vec{a} = (6\sqrt{t}\hat{j})$, donde $\vec{a}$ está en metros por segundo cuadrado y $t$ está en s. (a) Determine el vector velocidad de la partícula como función del tiempo. (b) Determine la posición de la partícula como función del tiempo.
    \vspace*{6cm}

    \item[\textbf{37.}] Lisa en su Lamborghini acelera a razón de $(3.00\hat{i} - 2.00\hat{j}) \text{ m/s}^2$, mientras que Jill en su Jaguar acelera a $(1.00\hat{i} + 3.00\hat{j}) \text{ m/s}^2$. Ambas parten desde el reposo en el origen de un sistema de coordenadas $xy$. Después de $5.00 \text{ s}$, (a) ¿cuál es la rapidez de Lisa respecto a Jill?, (b) ¿a qué distancia están?, y (c) ¿cuál es la aceleración de Lisa relativa a Jill?
    \vspace*{6cm}
\end{itemize}

\newpage

\begin{itemize}
    \item[\textbf{38.}] Un chico lanza una piedra horizontalmente desde la cima de un acantilado de altura $h$ hacia el océano. La piedra golpea el océano a una distancia $d$ desde la base del acantilado. En términos de $h$, $d$ y $g$, encuentre expresiones para el tiempo $t$ en que la piedra cae en el océano, (b) la rapidez inicial de la piedra, (c) la rapidez de la piedra justo antes de que llegue al océano y (d) la dirección de la velocidad de la piedra justo antes de que llegue al océano.
    \vspace*{7cm}

    \item[\textbf{39.}] ¿Por qué no es posible la siguiente situación? Albert Pujols hace un jonrón, de forma que la pelota libra la fila superior de las gradas, $24.0 \text{ m}$ de altura, situada a $130 \text{ m}$ de la base de \textit{home}. La pelota es golpeada a $41.7 \text{ m/s}$ en un ángulo de $35.0^\circ$ con la horizontal y la resistencia del aire es despreciable.
    \vspace*{6cm}

    \item[\textbf{40.}] Mientras algún metal fundido salpica, una gota vuela hacia el este con velocidad inicial $v_i$ a un ángulo $\theta_i$ sobre la horizontal y otra gota vuela hacia el oeste con la misma rapidez al mismo ángulo sobre la horizontal, como se muestra en la figura P4.40. En términos de $v_i$ y $\theta_i$, encuentre la distancia entre las gotas como función del tiempo.
    \vspace*{6cm}

    \item[\textbf{41.}] Un astronauta en la superficie de la Luna dispara un cañón para lanzar un paquete experimental, que deja el barril con movimiento horizontal. Suponga que la acelaración de caída libre en la Luna es un sexto de la propia de la Tierra. (a) ¿Cuál debe ser la rapidez de boquilla del paquete de modo que viaje completamente alrededor de la Luna y regrese a su ubicación original? (b) ¿Cuánto tarda este viaje alrededor de la Luna?
    \vspace*{6cm}
\end{itemize}

\newpage

\begin{itemize}
    \item[\textbf{42.}] Un péndulo con un cordón de longitud $r = 1.00 \text{ m}$ se balancea en un plano vertical (figura P4.42). Cuando el péndulo está en las dos posiciones horizontales $\theta = 90.0^\circ$ y $\theta = 270^\circ$, su rapidez es $5.00 \text{ m/s}$. (a) Encuentre la magnitud de la aceleración radial y (b) la aceleración tangencial para estas posiciones. (c) Dibuje diagramas vectoriales para determinar la dirección de la aceleración total para estas dos posiciones. (d) Calcule la magnitud y dirección de la aceleración total en estas dos posiciones.
    \vspace*{8cm}

    \item[\textbf{43.}] Un cañón de resorte se ubica en el borde de una mesa que está a $1.20 \text{ m}$ sobre el suelo. Una bola de acero se lanza desde el cañón con rapidez $v_i$ a $35.0^\circ$ sobre la horizontal. (a) Encuentre la posición horizontal de la bola como función de $v_i$ en el instante que aterriza. Esta función se escribe como $x(v_i)$. Evalúe $x$ para (b) $v_i = 0.100 \text{ m/s}$ y para (c) $v_i = 100 \text{ m/s}$. (d) Suponga que $v_i$ está cerca de cero, pero no es igual a cero. Muestre que un término en la respuesta al inciso (a) domina de modo que la función $x(v_i)$ se reduce a una forma más simple. (e) Si $v_i$ es muy grande, ¿cuál es la forma aproximada de $x(v_i)$? (f) Describa la forma global de la gráfica de la función $x(v_i)$.
    \vspace*{8cm}

    \item[\textbf{44.}] Se lanza un proyectil desde el punto $(x = 0, y = 0)$ con velocidad $(12.0\hat{i} + 49.0\hat{j})$ en $t = 0$. (a) Tabule la distancia del proyectil $|\vec{r}|$ desde el origen al final de cada segundo de allí en adelante, para $0 \le t \le 10 \text{ s}$. También puede ser útil tabular las coordenadas $x$ y $y$, y las componentes de velocidad $v_x$ y $v_y$. (b) Observe que la distancia del proyectil desde su punto de partida aumenta con el tiempo, llega a un máximo y comienza a disminuir. Pruebe que la distancia es un máximo cuando el vector de posición es perpendicular a la velocidad. \textit{Sugerencia:} Argumente que si $\vec{v}$ no es perpendicular a $\vec{r}$, entonces $|\vec{r}|$ debe aumentar o disminuir. (c) Determine la magnitud de la distancia máxima. (d) Explique su método para resolver el inciso (c).
    \vspace*{8cm}
\end{itemize}

\newpage

\begin{itemize}
    \item[\textbf{45.}] Un pescador emprende el viaje a contracorriente. Su pequeño bote, impulsado por un motor fuera de borda, viaja con rapidez constante $v$ en aguas tranquilas. El agua circula con rapidez constante $v_a$, menor. Recorre $2.00 \text{ km}$ a contracorriente cuando su hielera cae del bote. Se da cuenta de la falta de la hielera solo después de otros 15 minutos de ir a contracorriente. En ese punto, regresa río abajo, todo el tiempo viajando con la misma rapidez respecto al agua. Alcanza a la hielera justo cuando está próxima a la cascada en el punto de partida. ¿Con qué rapidez se mueven en las aguas del río? Resuelva este problema en dos formas. (a) Primero, use la Tierra como marco de referencia. Respecto de la Tierra, el bote viaja a contracorriente con rapidez $v - v_a$ y río abajo a $v + v_a$. (b) Una segunda solución mucho más simple y más elegante se obtiene al usar el agua como marco de referencia. Este planteamiento tiene importantes aplicaciones en problemas mucho más complicados; por ejemplo, en el cálculo del movimiento de cohetes y satélites y en el análisis de la dispersión de partículas subatómicas de objetivos de gran masa.
    \vspace*{7cm}

    \item[\textbf{46.}] Un jardinero lanza una pelota de béisbol a su receptor en un intento de sacar un corredor del cuadro. La pelota rebota una vez antes de llegar al receptor. Suponga que el ángulo al que una pelota que rebota deja el suelo es el mismo que el ángulo al que el jardinero la lanzó, como se muestra en la figura P4.46, pero la rapidez de la pelota después del rebote es un medio de la que tenía antes del rebote. (a) Suponga que la pelota siempre se lanza con la misma rapidez inicial. ¿A qué ángulo $\theta$ el jardinero debe lanzar la pelota para hacer que recorra la misma distancia $D$ con un rebote (trayectoria azul) que una lanzada hacia arriba a $45.0^\circ$ sin rebote (trayectoria verde)? (b) Determine la relación del intervalo de tiempo para el lanzamiento de un rebote y el tiempo de vuelo para el lanzamiento sin rebote.
    \vspace*{7cm}

    \item[\textbf{47.}] No se lastime; no golpee su mano contra nada. Dentro de estas limitaciones, describa cómo dar a la mano una gran aceleración. Calcule el orden de magnitud de esta aceleración, indicando las cantidades que mide o estima y sus valores.
    \vspace*{6cm}

    \item[\textbf{48.}] Usted está en la atracción de Piratas del Caribe en el Magic Kingdom en Disney World. Su bote se pasea por una batalla pirata, en la que los cañones de un barco y de un fuerte se disparan el uno al otro. Mientras que usted es consciente de que las salpicaduras en el agua no representan balas de cañón reales, comienza a preguntarse sobre tales batallas en los días de los piratas. Suponga que la fortaleza y la nave están separadas por $75.0 \text{ m}$. Ve que los cañones en el fuerte están dirigidos para que las balas de cañón sean disparadas horizontalmente desde una altura de $7.00 \text{ m}$ sobre el agua. (a) Usted se pregunta a qué velocidad se deben disparar para que peguen en el barco antes de caer en el agua. (b) Luego, usted piensa en el lodo que se debe acumular dentro del barril de un cañón. Este lodo debe frenar las balas de cañón. Piensa otra pregunta: si las balas de cañón se pueden disparar a solo $50.0\%$ de la velocidad encontrada antes, ¿es posible disparar hacia arriba en algún ángulo respecto horizontal para que alcancen la nave?
    \vspace*{7cm}
\end{itemize}

\newpage

\subsection*{PROBLEMAS DE DESAFÍO}

\begin{itemize}
    \item[\textbf{49.}] Un esquiador deja una rampa de salto con una velocidad de $10.0 \text{ m/s}$, a $15.0^\circ$ sobre la horizontal, como se muestra en la figura P4.49. La pendiente está inclinada a $\phi = 50.0^\circ$ y la resistencia del aire es despreciable. Encuentre (a) la distancia desde la rampa hasta donde aterriza el esquiador y (b) las componentes de velocidad justo antes de aterrizar. (c) Explique ¿cómo cree que afectan los resultados si se incluye resistencia del aire?
    \vspace*{8cm}

    \item[\textbf{50.}] Un proyectil es disparado hacia arriba en una pendiente (ángulo de inclinación $\phi$) con una velocidad inicial $v_i$ a un ángulo $\theta_i$ respecto a la horizontal ($\theta_i > \phi$) como se muestra en la figura P4.50. (a) Demuestre que el proyectil viaja una distancia $d$ hacia arriba de la pendiente, donde
    $$ d = \frac{2v_i^2 \cos\theta_i \operatorname{sen}(\theta_i - \phi)}{g \cos^2\phi} $$
    (b) ¿Para qué valor de $\theta_i$ es $d$ un máximo, y cuál es ese valor máximo?
    \vspace*{8cm}

    \item[\textbf{51.}] Dos nadadores, Chris y Sarah, parten desde el mismo punto en la orilla de una corriente ancha que circula con una rapidez $v$. Ambos se mueven con la misma rapidez $c$ (donde $c > v$) en relación con el agua. Chris nada corriente abajo una distancia $L$ y luego corriente arriba la misma distancia. Sarah nada de modo que su movimiento en relación con la Tierra es perpendicular a las orillas de la corriente. Ella nada la distancia $L$ y luego de vuelta la misma distancia, de modo que ambos nadadores regresan al punto de partida. En términos de $L, c$ y $v$, encuentre los intervalos de tiempo requeridos. (a) Para Chris de ida y vuelta y (b) para Sarah de ida y vuelta. (c) Explique cuál nadador regresa primero.
    \vspace*{8cm}
\end{itemize}

\newpage

\begin{itemize}
    \item[\textbf{52.}] En la sección \textbf{¿Qué pasaría si?} del ejemplo 4.5, se afirmó que el intervalo máximo de un esquiador se presenta para un ángulo de lanzamiento $\theta$ dado por
    $$ \theta = 45^\circ - \frac{\phi}{2} $$
    donde $\phi$ es el ángulo que la colina forma con la horizontal en la figura 4.15. Compruebe esta afirmación al derivar esta ecuación.
    \vspace*{7cm}

    \item[\textbf{53.}] Un cohete de fuegos artificiales explota a una altura $h$, en el pico de su trayectoria vertical. Lanza fragmentos ardientes en todas direcciones, pero todos con la misma rapidez $v$. Pedazos de metal solidificado caen al suelo sin resistencia del aire. Encuentre el ángulo más pequeño que hace la velocidad final de un fragmento que impacte, respecto a la horizontal.
    \vspace*{7cm}
\end{itemize}

\end{document}
